\documentclass[11pt]{article}
\usepackage[margin=1in]{geometry}
\usepackage{amsmath,amssymb,amsthm}
\usepackage[T1]{fontenc}
\usepackage{lmodern}
\usepackage[colorlinks=true,allcolors=blue]{hyperref}
\newtheorem{theorem}{Theorem}
\newtheorem{lemma}{Lemma}
\newtheorem{proposition}{Proposition}
\newcommand{\R}{\mathbb R}
\newcommand{\grad}{\operatorname{grad}}
\newcommand{\Ric}{\operatorname{Ric}}
\title{Radial orthogonality does not imply\\dual $1$-conformal flatness}
\author{Alec Kriebel\\\small Independent researcher\\
\small\href{https://orcid.org/0009-0001-9320-500X}{ORCID: 0009-0001-9320-500X}}
\date{22 September 2026\quad (version 1.0.0)}
\begin{document}
\maketitle
\begin{abstract}
We give a negative answer to Kurose's Question 3(e) in the 1998 collection
\emph{Open Problems in Affine Differential Geometry and Related Topics}.
For the product $S^2(1)\times\R$ with its product metric and Levi--Civita
connection, every radial orthogonal distribution in that question is
integrable, but the dual statistical structure is nowhere locally
$1$-conformally flat. Integrability follows from the first variation of
energy; nonflatness follows from two curvature evaluations. A conformal
metric and projective connection change also give an example with nowhere
zero cubic tensor. The same obstruction works in every dimension at least
three. The proof is independent of computation.
\end{abstract}

\section{The question and the counterexample}
A statistical manifold $(M,\nabla,h)$ consists here of a positive-definite
metric $h$ and a torsion-free connection $\nabla$ such that $\nabla h$ is
totally symmetric. The dual connection $\nabla^*$ is determined by
\[
Xh(Y,Z)=h(\nabla_XY,Z)+h(Y,\nabla_X^*Z).
\]
Fix $p\in M$ and a $\nabla$-convex neighborhood $U$ of $p$. For
$q\in U\setminus\{p\}$, let $\gamma_q:[0,1]\to U$ be the affinely
parametrized $\nabla$-geodesic from $p$ to $q$. Set
\begin{equation}\label{radial}
D^p_q=\{w\in T_qM:h_q(\dot\gamma_q(1),w)=0\}.
\end{equation}
Only the velocity line matters, so a regular reparametrization does not
change this distribution.

Question 3(e) of \cite[p.~126]{FMU} asks whether integrability of every
distribution \eqref{radial} implies that $(M,\nabla^*,h)$ is
$1$-conformally flat. This is the printed question indexed as
AMR-059-0011 (6000011). It excludes the center $p$ and does not exclude
Levi--Civita connections; item 3(a) explicitly treats those as statistical
structures.

Following \cite[p.~428]{Kurose}, a $1$-conformal change is
\begin{equation}\label{change}
\widehat h=e^\varphi h,\qquad
\widehat\nabla_XY=\nabla_XY-h(X,Y)\grad_h\varphi.
\end{equation}
The structure is $1$-conformally flat if such a change makes
$\widehat\nabla$ flat locally at every point. This is a condition on the
statistical connection, distinct from ordinary conformal flatness of $h$.
We use $R(X,Y)Z=\nabla_X\nabla_YZ-\nabla_Y\nabla_XZ-\nabla_{[X,Y]}Z$.

\begin{theorem}\label{main}
Let $M=S^2(1)\times\R$, $h=g_{S^2(1)}+dt^2$, and
$\nabla=\nabla^h$. Every distribution \eqref{radial} is integrable,
for every center and every convex neighborhood on which it is defined.
Nevertheless $(M,\nabla^*,h)$ is not $1$-conformally flat on any nonempty
open subset. The same conclusions hold for
$S^2(1)\times\R^{n-2}$ in every dimension $n\geq3$.
\end{theorem}

\section{An intrinsic proof}
\begin{lemma}\label{energy}
For any Riemannian manifold with its Levi--Civita connection, every
distribution \eqref{radial} is integrable.
\end{lemma}
\begin{proof}
On a convex neighborhood define the smooth endpoint energy
\[
E_p(q)=\frac12\int_0^1 h(\dot\gamma_q(u),\dot\gamma_q(u))\,du.
\]
For an endpoint variation $q(s)$ with $q'(0)=w$, put
$F(s,u)=\gamma_{q(s)}(u)$, $T=\partial_uF$, and $V=\partial_sF$.
Metric compatibility, torsion-freeness, and the geodesic equation yield
\begin{align*}
dE_p|_q(w)
 &=\int_0^1h(\nabla_VT,T)\,du
 =\int_0^1h(\nabla_TV,T)\,du\\
 &=\bigl[h(V,T)\bigr]_0^1-\int_0^1h(V,\nabla_TT)\,du
 =h_q(w,\dot\gamma_q(1)).
\end{align*}
The initial endpoint is fixed, so $V(s,0)=0$. For $q\ne p$ the terminal
velocity is nonzero, and therefore
$dE_p(\dot\gamma_q(1))=h(\dot\gamma_q(1),\dot\gamma_q(1))>0$.
Thus $D^p=\ker dE_p$ on $U\setminus\{p\}$, and its integral
hypersurfaces are the regular level sets of $E_p$. On sufficiently small
normal neighborhoods these are the geodesic spheres, since
$E_p(q)=\tfrac12 d_h(p,q)^2$. The variation argument itself applies
on any convex neighborhood in the question.
\end{proof}

\begin{lemma}\label{obstruction}
For a Levi--Civita statistical structure, local $1$-conformal flatness
requires an endomorphism $A$ at each point satisfying
\begin{equation}\label{necessary}
R(X,Y)Z=h(Y,Z)A(X)-h(X,Z)A(Y).
\end{equation}
\end{lemma}
\begin{proof}
Write $V=\grad_h\varphi$ and $B(X,Y)=-h(X,Y)V$ for the difference
tensor in \eqref{change}. The curvature transformation law gives
\begin{align*}
\widehat R(X,Y)Z
 &=R(X,Y)Z+(\nabla_XB)(Y,Z)-(\nabla_YB)(X,Z)\\
 &\hspace{1em}+B(X,B(Y,Z))-B(Y,B(X,Z))\\
 &=R(X,Y)Z-h(Y,Z)A(X)+h(X,Z)A(Y),
\end{align*}
where $A(X)=\nabla_XV-h(X,V)V$. Indeed,
$(\nabla_XB)(Y,Z)=-h(Y,Z)\nabla_XV$ and
$B(X,B(Y,Z))=h(Y,Z)h(X,V)V$.
Setting $\widehat R=0$ proves \eqref{necessary}.
\end{proof}

\begin{proof}[Proof of Theorem~\ref{main}]
Metric compatibility gives $\nabla h=0$, so the structure is statistical
and $\nabla^*=\nabla$. Lemma~\ref{energy} gives the required integrability.
At any point choose orthonormal vectors $e_1,e_2$ tangent to the unit
sphere and $e_3$ tangent to the line. Product curvature gives
\[
R(e_1,e_2)e_2=e_1,\qquad R(e_1,e_3)e_3=0.
\]
If \eqref{necessary} held, these two equations would give respectively
$A(e_1)=e_1$ and $A(e_1)=0$, a contradiction. This excludes even an
arbitrary endomorphism $A$, before imposing its differential form in
Lemma~\ref{obstruction}. The contradiction holds at every point.
Additional Euclidean factors leave these evaluations unchanged.
\end{proof}

For comparison, Kurose's criterion \cite[Proposition~1]{Kurose} identifies
$1$-conformal flatness with projective flatness of the dual connection
with symmetric Ricci tensor. Here $\Ric=\operatorname{diag}(1,1,0)$
in the chosen frame. The projective Weyl tensor is
\[
W(X,Y)Z=R(X,Y)Z-\tfrac12\{\Ric(Y,Z)X-\Ric(X,Z)Y\},
\qquad W(e_3,e_1)e_1=-\tfrac12e_3\ne0.
\]
This is an independent obstruction. Neither this criterion nor a
classification theorem is needed for the proof above.

\section{A non-self-dual example}
The failure is not confined to zero cubic tensor. Let
$h_0=g_{S^2(1)}+dt^2$, $\nabla^0=\nabla^{h_0}$, and define
\begin{equation}\label{extension}
h_1=e^t h_0,\qquad
\nabla^1_XY=\nabla^0_XY+dt(X)Y+dt(Y)X.
\end{equation}
This is a standard projective change of connection, combined with a
conformal change of metric. Such mechanisms belong to the theory of
divisible cubic forms; see, for example, \cite[\S3.2]{Ueno}. We use
them only to remove the possible restriction to self-dual examples.

\begin{proposition}
The structure $(M,\nabla^1,h_1)$ has nowhere zero cubic tensor and
integrable radial distributions for every center. Its dual is nowhere
locally $1$-conformally flat.
\end{proposition}
\begin{proof}
The connection is torsion-free, and direct differentiation gives
\begin{equation}\label{cubic}
(\nabla^1_Xh_1)(Y,Z)=-dt(X)h_1(Y,Z)-dt(Y)h_1(X,Z)-dt(Z)h_1(X,Y).
\end{equation}
This is symmetric and has value $-3e^t$ on
$(\partial_t,\partial_t,\partial_t)$. The connections $\nabla^1$
and $\nabla^0$ have the same unparametrized geodesics:
\[
\nabla^1_{\dot\gamma}\dot\gamma
=\nabla^0_{\dot\gamma}\dot\gamma+2dt(\dot\gamma)\dot\gamma.
\]
The additional acceleration is tangent to the curve and is removed by
a regular reparametrization. More explicitly, if $u$ is a
$\nabla^0$-affine parameter then a $\nabla^1$-affine parameter $s$
satisfies $ds/du=C e^{2t(\gamma(u))}$ with $C>0$. This works on each
compact geodesic segment. It identifies the unique connecting paths in
the corresponding convex neighborhoods. Since $h_1$ and $h_0$ have the
same orthogonality relation, their radial distributions agree under
this identification, so Lemma~\ref{energy} applies.

The defining duality identity gives
\[
(\nabla^1)^*_XY=\nabla^0_XY-h_0(X,Y)\partial_t.
\]
For any local function $\psi$, its $1$-conformal change is
\begin{align*}
\widehat\nabla_XY
 &=(\nabla^1)^*_XY-h_1(X,Y)\grad_{h_1}\psi\\
 &=\nabla^0_XY-h_0(X,Y)\grad_{h_0}(t+\psi).
\end{align*}
If this were flat, it would be a flat $1$-conformal change of
$(\nabla^0,h_0)$, contrary to Theorem~\ref{main}.
\end{proof}

\section*{Verification, attribution, and scope}
The obstruction is elementary and the ingredients are classical.
The contribution asserted here is the explicit negative answer to the
printed Question 3(e), together with a short direct verification.
A focused literature audit found no earlier explicit resolution, but
does not establish historical priority. In particular, the live problem
database's annotations could not be checked; its identifier is only a
cross-reference. We do not claim that the Gauss lemma, the curvature
criterion, or the construction in \eqref{extension} is new.

The accompanying standard-library Python checker solves the complete
pointwise linear system \eqref{necessary} in exact rational arithmetic.
A separately implemented SymPy checker derives the Christoffel symbols,
curvature, transformation law, and cubic and dual tensors from
$h_0=d\theta^2+\sin^2\theta\,d\phi^2+dt^2$. These checks supplement
the proof; they do not replace its quantified first-variation argument
or constitute a proof-assistant formalization. No dimension-two claim
or global foliation across cut loci is needed.

\paragraph{AI disclosure.}
The candidate argument was supplied by the author as AI-generated work.
This manuscript, supporting computations, source audit, and separate
adversarial AI reviews were prepared with OpenAI Codex. This is an
unrefereed preprint; AI reviews are not external human peer review.

\begin{thebibliography}{9}
\bibitem{FMU}
H.~Furuhata, H.~Matsuzoe and H.~Urakawa,
Open Problems in Affine Differential Geometry and Related Topics,
\emph{Interdisciplinary Information Sciences} \textbf{4} (1998), 125--127.
\href{https://doi.org/10.4036/iis.1998.125}{doi:10.4036/iis.1998.125}.

\bibitem{Kurose}
T.~Kurose, On the divergences of 1-conformally flat statistical manifolds,
\emph{Tohoku Mathematical Journal} \textbf{46} (1994), 427--433.
\href{https://doi.org/10.2748/tmj/1178225722}{doi:10.2748/tmj/1178225722}.

\bibitem{Ueno}
R.~Ueno, Geodesic Connectedness on Statistical Manifolds with Divisible
Cubic Forms, arXiv:2503.10024v2 (2025).
\url{https://arxiv.org/abs/2503.10024v2}.
\end{thebibliography}
\end{document}
